\scanpage{354}
\mspage{478}

that $N_0$ and $N_1$ be everywhere
\unsure{non-homothetic|nonhomothetic}.

Then $\widetilde\varepsilon_0(\mathbb G_a)$,
$\widetilde\varepsilon_1(\mathbb G_a)$,
$\widetilde\varepsilon_0\in\Gamma
\widetilde\varepsilon_0(\mathbb G_a)$ define
a canonical \deleted{\unsure{splitting}}
\inserted{pinning} of $PG$ hence of $G$ \ldots
\inserted{We will then have, for this pinning}
\deleted{The situation is}
\[
(26)\qquad
\rho=\begin{pmatrix}0&u\\v&w\end{pmatrix}
\qquad
\sigma=\begin{pmatrix}0&u'\\v'&w'\end{pmatrix}
\]
with $u,v,u',v'\in\Gamma(\mathbb G_m)$,
$w,w'\in\Gamma\mathbb G_a$,
\deleted{$\unsure{\rho}$}

and we will have
\[
(27)\quad
\left\{
\begin{aligned}
-\sigma^{-1}\rho
&=\begin{pmatrix}
-\dfrac{v}{v'}&
\dfrac{uw'-wu'}{u'v'}\\
0&-\dfrac{u}{u'}
\end{pmatrix}
=\varepsilon_0
=\begin{pmatrix}1&\unsure{?}\\0&1\end{pmatrix}\\[1ex]
-\rho\sigma^{-1}
&=\begin{pmatrix}
-\dfrac{u}{\unsure{w'|u'}}&0\\
\dfrac{vw'\unsure{-|+}wv'}{\unsure{w'|u'}v'}&-\dfrac{v}{v'}
\end{pmatrix}
=\varepsilon_1
\overset{\inserted{\unsure{\mathrm{déf.}}}}{=}
\begin{pmatrix}1&0\\\unsure{1}&1\end{pmatrix}
\end{aligned}
\right.
\]
\margin{by definition of our pinning}

so we will need
\[
(28)\qquad
u/u'=-1,\quad v/v'=-1,\qquad
\text{i.e.}\quad u'=-u,\ v'=-v
\]
\deleted{$(29)\quad\dfrac{vw'\unsure{+|-}wv}{uv}$}
so
\[
-\rho\sigma^{-1}
=\begin{pmatrix}1&0\\\dfrac{1}{u}(w'+w)&1\end{pmatrix}
\]
and we need
\[
(29)\qquad w'+w=u
\qquad\text{i.e.}\qquad w'=w-u
\]
in other words \unsure{combining} $(26)(27)$
\[
(30)\qquad \rho+\sigma=uP
\quad\text{where}\quad
P=\begin{pmatrix}0&0\\0&1\end{pmatrix}
\]
is the elementary projector associated with the pinning under consideration.
\margin{which determines $\sigma$ in terms of $\rho$,
\deleted{and} the \unsure{coefficient} $u$ of $\rho$, and $P$}

And we will have
\[
(31)\qquad
-\sigma^{-1}\rho=\varepsilon_0
=\begin{pmatrix}1&u/v\\0&1\end{pmatrix},
\qquad
N_0=\begin{pmatrix}0&u/v\\0&\unsure{0|1}\end{pmatrix}
\]

\mspage{353}

\scanpage{355}

We would like to give intrinsic ``numerical criteria'' for $\rho,\sigma$ to have the preceding properties\ldots

\uline{Proposition} Let $\rho,\sigma$ be two \inserted{\deleted{\unsure{regular}}} sections of $G$,

\deleted{Equivalent conditions} \inserted{\unsure{Put}}
\[
\unsure{\varepsilon}\qquad \varepsilon_0=-\sigma^{-1}\rho,\qquad \varepsilon_1=-\rho\sigma^{-1}
\]
so that we thus have
\[
\varepsilon_1=\sigma(\varepsilon_0)=\rho(\varepsilon_0).
\]

Equivalent conditions

$a)$ \inserted{$1^\circ)$} $\varepsilon_0$ \deleted{and $\varepsilon_1$} \deleted{\unsure{are}} \inserted{is} \unsure{regular|regular} \unsure{unipotent|unipotent} \deleted{\unsure{everywhere}} \inserted{$(\Leftrightarrow\varepsilon_1$} regular unipotent) and \inserted{$2^\circ)$} $\varepsilon_0$ and $\varepsilon_1$ \deleted{everywhere} \deleted{\unsure{not equivalent, i.e.}} independent \inserted{in each fiber}, --- i.e. putting
\[
\varepsilon_0=1+N_0,\qquad \varepsilon_1=1+N_1
\]
$N_0$ $(\Leftrightarrow N_1)$ is regular nilpotent, and $N_0$ and $N_1$ are everywhere nonhomothetic.

$b)$ We have
\[
\begin{array}{rl}
\deleted{\unsure{(32)}}&\\[-1.2em]
(32)&\left\{\begin{array}{l}
\left.\begin{array}{ll}
a)&\det\sigma=\det\rho\\
b)&\operatorname{Tr}\sigma^{-1}\rho=-2
\end{array}\right\}\\
c)\quad u\overset{\mathrm{dfn}}{=}\operatorname{Tr}(\sigma+\rho)\text{ is invertible}
\end{array}\right.
\end{array}
\]
\deleted{$\unsure{(33)}$} \deleted{$\unsure{(34)}$}
\margin{these two relations express $\det(\varepsilon_0)=1$, $\operatorname{Tr}\varepsilon_0=2$ i.e. $\varepsilon_0$ is \deleted{(strictly)} unipotent}
\margin{this expresses that $\varepsilon_0$ is regular and independent of $\varepsilon_1$, cf. loc. cit.}
\margin{cf. proposition page $449'$}

$c)$ We have \deleted{$\det\sigma=\det\rho$}
\[
(33)\quad\left\{\begin{array}{ll}
a)&\det\sigma=\det\rho\\
b)&\det(\sigma+\rho)=0\\
c)&\operatorname{Tr}(\sigma+\rho)\text{ invertible}
\end{array}\right.
\]

$d)$ We have $\det\sigma=\det\rho$, and there exist an elementary projector $P$ and an invertible scalar $u\in\Gamma\mathbb G_{\mathrm m}$, such that we have
\[
(34)\qquad \sigma+\rho=uP
\]
\margin{\uline{N.B.} $P$ and $u$ are uniquely determined, $u=\operatorname{Tr}(\sigma+\rho)$, $P=\frac1u(\sigma+\rho)$}

\mspage{354}

\scanpage{356}
\mspage{479}

\deleted{$e)$ (If $M=\underline{\operatorname{End}}(F)$) $\exists$ a basis of $F$, such that with respect to this basis we have}

$e)$ $\exists$ isomorphism $\mathrm{G}\ell(2)\overset{\varphi}{\xrightarrow{\sim}}G$, \inserted{\deleted{(unique)}} such that we have
\[
(35)\quad\left\{\begin{aligned}
\varphi^{-1}(\varepsilon_1)&=\begin{pmatrix}1&0\\1&1\end{pmatrix}\\
\varphi^{-1}(\varepsilon_0)&=\begin{pmatrix}1&\deleted{\unsure{1}}\unsure{?}\\0&1\end{pmatrix}
\end{aligned}\right.
\]
\margin{\uline{N.B.} this isomorphism is unique}

$f)$ $\exists$ isomorphism $\mathrm{G}\ell(2)\simeq G$ such that we have
\[
(36)\quad\left\{\begin{gathered}
\varphi^{-1}(\rho)=\begin{pmatrix}0&u\\v&w\end{pmatrix}\qquad
\varphi^{-1}(\sigma)=\begin{pmatrix}0&\unsure{w'|u'}\\v'&w'\end{pmatrix}\\
\text{with}\qquad u+u'=v+v'=0,\qquad w+w'=u
\end{gathered}\right.
\]
\margin{\uline{N.B.} We will have $P=\begin{pmatrix}0&0\\0&1\end{pmatrix}$}
\margin{this $\varphi$ is again unique; We will have
\[
(37)\quad\left\{\begin{aligned}
u&=\operatorname{Tr}(\sigma+\rho)\;(=-u')\\
w&=\operatorname{Tr}\rho\quad(w'=\operatorname{Tr}\sigma)\\
-uv&=\det\rho=\det\sigma
\end{aligned}\right.
\]
which describes $u,\unsure{w},\unsure{v}$ in terms of $\rho$, hence also $u',v',w'\ldots$}

\deleted{$\unsure{(37)}$}

We will have in (35), with
\[
\unsure{(38)}\qquad\varphi^{-1}(\varepsilon_0)=\begin{pmatrix}1&u/v\\0&1\end{pmatrix},
\]
and we will have
\[
(39)\qquad\sigma+\rho=uP,\quad\text{with }\varphi^{-1}(P)=\begin{pmatrix}0&0\\0&1\end{pmatrix}.
\]

Easy proof. The equivalence $a)\Leftrightarrow b)$ \unsure{seen|seen} in proposition page $449'$. For $b)\Leftrightarrow c)$, we write
\[
\begin{aligned}
\unsure{\det}(\sigma+\rho)
&=\det\{\sigma(1+\sigma^{-1}\rho)\}
=\det\sigma\;\underbrace{\det(1+\sigma^{-1}\rho)}_{1+\operatorname{Tr}\sigma^{-1}\rho+\underbrace{\det\sigma^{-1}\rho}_{(\det\sigma)^{-1}\det\rho}}\\
&=\\
&=\det\sigma+\det\rho+(\det\sigma)\operatorname{Tr}(\sigma^{-1}\rho)
\end{aligned}
\]
so by virtue of the hypothesis
\[
\det\sigma=\det\rho
\]
common to $b)$ and $c)$, \unsure{we will have}
\[
(\unsure{*|6})\qquad\det(\sigma+\rho)=\det\sigma\,(2+\operatorname{Tr}(\sigma^{-1}\rho))
\]
so the first \unsure{side|side} is zero iff $\operatorname{Tr}\sigma^{-1}\rho=-2$, which was to be proved.

\mspage{355}

\scanpage{357}

Putting $Q=\sigma+\rho$, the equivalence $b\Leftrightarrow c$ comes from the equivalence (valid for every section $Q$ of $M$)
\[
(40)\quad (\det Q=0,\operatorname{Tr}Q\text{ invertible})\Longleftrightarrow
\left(\begin{gathered}\exists\text{ elementary projector }P,\text{ and }u\in\Gamma\mathbb G_{\mathrm m},\\
\text{such that }Q=uP\end{gathered}\right)
\]
which follows easily from the results of the preceding n\textsuperscript{o} (it is a corollary of the proposition page $475'$, which I forgot to make explicit.

The equivalence with $e)$, $f)$ is also immediate, \deleted{\unsure{it is in|it is a}} and merely summarizes certain preceding calculations.

\margin{\uline{Corollary} Under these conditions $1,N_0,N_1,P$ form a \uline{basis} of $M$, and $\rho$ and $\sigma$ are independent linear combinations of \deleted{$\unsure{N_0}$} $N_0,N_1,P$
\[
\left\{\begin{aligned}
\rho&=v(N_0+N_1)+wP\\
\sigma&=-v(N_0+N_1)+(u-w)P
\end{aligned}\right.
\]
$\rho,\sigma$ form an \unsure{obviously|obvious} basis of the \unsure{sub-Module|submodule} of $M$ generated by $(N_0+N_1)$ and $P$, hence
\[
\rho,\sigma,N_0,1\qquad\text{and}\qquad\rho,\sigma,N_1,1
\]
also form bases of $M$, and in them we can replace either $\rho$, or $\sigma$ by $P$; i.e. in $N_0,N_1,P,1$ we can \unsure{replace} $N_0$, or $N_1$ by $\rho$ or $\sigma$.}

\figure{\begin{tikzpicture}[x=1cm,y=1cm]
\draw[->] (0,1.5) .. controls (-0.15,0.75) and (0.45,0.6) .. (1.05,0.85) -- (-0.05,-0.9);
\end{tikzpicture}}

Our point of view here is that we start from $\widetilde\rho,\widetilde\sigma_{\unsure{1}}$, which we want to describe --- the pinning \unsure{is} indeed already defined in terms of $\widetilde\rho,\widetilde\sigma$, but there remains the question of normalizing the liftings of $\widetilde\rho,\widetilde\sigma$ to $\rho,\sigma$. We have already made a first restriction by the formulas
\[
\varepsilon_1=-\sigma^{-1}\rho
\]
by eliminating the scalar \inserted{\uline{factor}} $\lambda_0$ --- we can \unsure{normalize|simplify} further by multiplying $\rho,\sigma$ by the factor $u^{-1}$, so as to obtain the simple expression.
\[
(40)\qquad\operatorname{Tr}(\rho+\sigma)=1\quad\text{i.e. }\operatorname{Tr}\rho+\operatorname{Tr}\sigma=1\quad\text{i.e. }u=1,
\]
which thus means that $\rho,\sigma$ are of the

\mspage{356}

\scanpage{358}
\mspage{480}

form
\[
(41)\qquad \rho=\begin{pmatrix}0&1\\v&w\end{pmatrix},\qquad
\sigma=\begin{pmatrix}0&-1\\v'&w'\end{pmatrix}
\]
\[
v+v'=0,\quad w+w'=1\qquad\text{i.e. }v'=-v,\quad w'=1-w,
\]
so
\[
(42)\qquad\varepsilon_0=\begin{pmatrix}1&1/v\\0&1\end{pmatrix},
\]
where we have identified, by $\varphi$, $\mathrm{G}\ell(2)_S$ and $G$.

\deleted{We \unsure{will note} that $\varepsilon_1$ is \unsure{entirely determined}}
\inserted{Is. there uniqueness of $\rho,\sigma$ in terms of $\widetilde\rho,\widetilde\sigma$,}
\deleted{\unsure{by}} by the conditions $\varepsilon_1=-\rho\sigma^{-1}$ (i.e. $\varepsilon_0=-\sigma\rho^{-1}$) and (40)?
I would first like to check carefully an assertion made in passing a certain number of times, namely that the \deleted{map} \inserted{morphism}
\[
G\longrightarrow PG\simeq G/\mathbb G_{\mathrm m}
\]
induces a \uline{bijection}

(43)\quad
\diagram{\begin{tikzcd}[column sep=large,ampersand replacement=\&]
\begin{array}{c}\text{unipotents }\deleted{\unsure{strict}}\\\text{regular of}\\G\end{array}
\arrow[r] \&
\begin{array}{c}\text{unipotents}\\\text{\uline{regular}}\\\text{of }PG\end{array}
\end{tikzcd}},

(at least locally \deleted{\unsure{fppf}} \inserted{\unsure{fppf}} \deleted{\unsure{Zariski}})

the surjectivity of this map could here serve as a \uline{definition} for the regular unipotents of $PG$. So it is a matter of testing \uline{injectivity}, which also means \deleted{that} this: if $\varepsilon,\varepsilon'$ are two regular unipotent sections of $M$, such that \inserted{$\widetilde\varepsilon=\widetilde\varepsilon'$ i.e.} $\varepsilon'=\lambda\varepsilon$ $(\lambda\in\Gamma\mathbb G_{\mathrm m})$, then $\varepsilon=\varepsilon'$, \inserted{i.e. $\lambda=1$.} Taking the determinants of the two sides, we will have
\[
(44)\qquad\lambda^2=1,
\]
\margin{So, we will say that a section of $PG$ is \unsure{regular} unipotent if it is locally fppf the image of a regular unipotent section of $G$.}

\mspage{357}